\documentclass[11pt,a4paper]{article}
\usepackage[margin=19mm,headheight=15pt]{geometry}
\usepackage{fontspec}
\setmainfont{Noto Serif}
\usepackage{amsmath,amssymb,mathtools}
\usepackage{xcolor}
\usepackage{fancyhdr}
\usepackage{enumitem}
\usepackage{microtype}
\setlength{\parindent}{0pt}
\setlength{\parskip}{6pt}
\definecolor{navy}{HTML}{183153}
\definecolor{lightgray}{HTML}{F3F5F7}
\pagestyle{fancy}\fancyhf{}
\lhead{\small ME 5309 -- Mid-Semester Examination}
\rhead{\small ME 5309}
\cfoot{\thepage}
\begin{document}
\begin{center}
{\Large\bfseries ME 5309 -- Mid-Semester Examination}\\[3pt]
{\large Answer Booklet}\\[8pt]
\end{center}
\textbf{Instructions:} Attempt any three questions. Each question carries 10 marks. Show working where appropriate.
\section*{Question 1 \hfill [10 marks]}
\subsection*{Q1(a) \hfill [4 marks]}

Pixel pitch = 0.005 mm, sensor width = 1280x0.005 = 6.4 mm. how to calculate the belt field of view.

\subsection*{Q1(b) \hfill [3 marks]}

\[\begin{aligned}Sensitivity = 10 mV/N\end{aligned}\]
The ADC has 12 bits, so there are 4096 levels.

\subsection*{Q1(c) \hfill [3 marks]}

\[\begin{aligned}F = e^(-2\times0.1) = 0.818731\end{aligned}\]
I used G = (1-F)/(-2) = -0.0906346.
So:\[\begin{aligned}x_{k+1} = 0.818731x_k - 0.0906346u_k\end{aligned}\]

\section*{Question 2 \hfill [10 marks]}
\subsection*{Q2(a) \hfill [2 marks]}

\vspace{2mm}

\subsection*{Q2(b) \hfill [2 marks]}

\[\begin{aligned}A = [[0,1], [-mg\ell sin(60^{\circ})/J, -b/J]] = [[0,1], [-106.2, -2.5]],\\ B = [[0],[50]]\end{aligned}\]

\subsection*{Q2(c) \hfill [4 marks]}

At rest, θ̇ = θ̈ = 0, so the torque balance is Kp(θd-θss) = mgℓ sin θss.
With Kp = 50, θd = π/2 and mgℓ = 2.4525, solve 50(π/2-θss) = 2.4525 sin θss.
This gives θss approximately 1.52181 rad = 87.193°, so the error is 90°-87.193° = 2.807°; it is not zero.\par\smallskip Rearranging the steady-state equation gives $50(\pi/2-\theta_{ss})=2.4525\sin\theta_{ss}$. I used radians in the equation. Solving numerically gives $\theta_{ss}\approx1.522\,$rad $\approx87.19^\circ$, so the gap from $90^\circ$ is about $2.81^\circ$.

\subsection*{Q2(d) \hfill [2 marks]}

Gravity compensation makes the steady-state error zero.
It does not change the closed-loop poles because it only adds a constant torque.

\section*{Question 3 \hfill [10 marks]}
\subsection*{Q3(a) \hfill [2 marks]}

The estimate is noisy and can fluctuate from sample to sample.

\subsection*{Q3(b) \hfill [2 marks]}

\[\begin{aligned}y[k]=0.8818y[k-1]+0.0591x_k+0.0591x[k-1], DC gain 1\end{aligned}\]

\subsection*{Q3(c) \hfill [4 marks]}

With L=[ℓ₁,ℓ₂]ᵀ and C=[1,0], A-LC=[[-ℓ₁,1], [-ℓ₂,-2]].
Its characteristic polynomial is det(sI-(A-LC)) = (s+ℓ₁)(s+2)+ℓ₂ = s²+(ℓ₁+2)s+(2ℓ₁+ℓ₂).
Desired poles at -10 give (s+10)²=s²+20s+100.
Hence ℓ₁+2=20 gives ℓ₁=18, and 2ℓ₁+ℓ₂=100 gives ℓ₂=64.\par\smallskip Expanding $(s+\ell_1)(s+2)+\ell_2$ gives $s^2+(\ell_1+2)s+(2\ell_1+\ell_2)$. The desired polynomial is $(s+10)^2=s^2+20s+100$. Matching the $s$ coefficient gives $\ell_1=18$; matching the constant term then gives $36+\ell_2=100$, hence $\ell_2=64$.

\subsection*{Q3(d) \hfill [2 marks]}

An observer uses the plant model and known control input to estimate unmeasured states such as velocity, and its poles can be designed.
A simple low-pass filter cannot exploit the state-space model in this way.
A filter may be safer when the model is poorly known or unreliable, because it relies less on model accuracy.\par\smallskip The observer uses the plant model and the measured output to reconstruct the unmeasured state. Its convergence depends on the model and the observer poles. A low-pass filter is simpler, but generally introduces phase lag and does not explicitly use the plant dynamics.

\section*{Question 4 \hfill [10 marks]}
\subsection*{Q4(a) \hfill [2 marks]}

Poles are -2±2j and Nff=8.

\subsection*{Q4(b) \hfill [2 marks]}

The nominal controller has unity DC gain, so θss=1 and the tracking error is zero even after the actuator gain changes.

\subsection*{Q4(c) \hfill [4 marks]}

\[\begin{aligned}k_1=18, k_2=6, K_I=32\end{aligned}\]

\subsection*{Q4(d) \hfill [2 marks]}

Integral action removes steady-state error, although this only holds for the nominal values of b and c.
Saturation can cause windup, which should be handled by anti-windup.

\end{document}